\section{Introduction}
An experiment does not acquire a manifold. It acquires reports through preparations and interventions, each of which has a causal position and a physical cost. The future tests executable after those reports determine which distinctions matter. Only after this experiment and its tests have been fixed does a geometric question arise: how much of the resulting predictive information can be retained and used with finite resources?

The central difficulty is not the definition of a predictive quotient. It is the passage from a quotient and an acquired law to an algorithm that uses the permitted interface, and to a converse that applies to all competing algorithms. Three obstacles occur together. Selection changes the acquired law and must pay for rejected observations. An error harmless for an ordinary report can be amplified by a rare stopping event. Finally, a label used to route a live protocol need not be unavailable as a terminal output. A geometric argument that omits any of these features can give an incorrect resource curve.

We address these obstacles by comparing \emph{unnormalized task measures}. Given a genuinely prepared joint report--audit law $K_0(dx,dz)$, let $\nu$ be its report law and derive
\[
 f(x)=\E_0[Z\mid x],\qquad w(x)=\E_0[\|Z-f(x)\|^2\mid x].
\]
Neither quantity is prescribed as an information metric. For a dictionary $C$ the task distortion is $w+\dist(f,C)^2$. A probability measure selected at expected paid-attempt budget $b$ is subject to a density cap, leading to the profile
\begin{equation}\label{eq:intro-profile}
 \Q_m(b)=\min_{|C|\le m,\ 0\le h\le b,\ \int h\,d\nu=1}
       \int [w+\dist(f,C)^2]h\,d\nu.
\end{equation}
This is a marked population trimmed-quantization functional. Trimming and finite dictionaries have a classical statistical history \cite{CGM97,GG99}; the claim here is not the invention of that functional. The problem is to prove when its geometry describes a paid, causal, unknown-model experiment.

Theorem~\ref{thm:transfer} is the central result. One actual operating model supplies a converse against \emph{every} legal strategy. A separate, parameter-independent acquisition program supplies an upper bound for \emph{every} operating model and every surviving history. We compare only the current productive report--audit row while retaining the actual pre-call occupation of the competitor. This yields an accepted comparison measure, not a supposed execution under the comparison experiment. Its mass, density cap and task loss can consequently be controlled without transferring a stopping time to an experiment in which the rule may never terminate.

The new feature is the simultaneous treatment of order and defect. Relative order of nonnegative task measures changes constants and effective budgets. Absolute defects are integrated against actual pre-call occupation. In particular, when there is relative order but no absolute defect, fixed nonzero calibration distortion does not impose an additive risk floor. Corollary~\ref{cor:uniform} gives matching resource curves at all scales. This differs from an estimate whose error grows as acquisition count times an absolute one-call discrepancy. Both regimes are consequences of the same theorem, rather than unrelated approximation results.

The order has a raw verification principle. The positive matrix measure
\[
 J(D)=\int_{D\times H}(1,z)(1,z)^\mathsf{T}\,P(dy,dz)
\]
contains exactly the localized affine-square tests. Proposition~\ref{prop:matrix-order} identifies the order of these tests with the order of $J$. Theorem~\ref{thm:singular-order} expresses it on the measurable range of $J$, including variable rank and zero-mass reports. It does not invert a vanishing eigenvalue uniformly. The construction is task-derived; it neither assumes a Fisher matrix nor identifies the matrix rank with a memory exponent.

There are two further structural conclusions. Theorem~\ref{thm:quotient} constructs the Borel predictive quotient and a measurable section from compact exhaustion and continuous response coordinates. It replaces an assumed Borel quotient realization by verifiable sufficient conditions; it does not assert that every measurable equivalence relation has such a realization. Lemma~\ref{lem:roles} gives separate exact packing rules for event and state stopping. The distinction persists under the resource-certified simulator composition of Proposition~\ref{prop:composition}.

The geometric consequences depend on the actual small-ball mass of the prepared law and on global covering. They permit singular location measures, a variable-rank task carrier and a positive acquired zero-loss atom. Two raw realizations verify the same general theorem: finite classical controlled channels and finite-dimensional noncommuting quantum instruments. Their targets are post-command audits. Neither is an estimation theorem for an untouched pre-intervention state. The full preceding theory, including nonintervening directional Gaussian sensing, absolute-calibration occupation bounds, finite-description examples and stopped-flow realization, is supplied as the complete technical companion \cite{Companion}. It is preserved, not relabelled as new mathematics.

\subsection*{Relation to experiment comparison and predictive representations}
Classical experiment comparison \cite{Blackwell} and predictive-state minimality \cite{SC} provide important antecedents. They do not by themselves give the particular paid stopping profile, a finite-state strategy lift, or a separate simulator clock and workspace account. Our theorem uses a reference/effort witness and a common task-order construction. It is not a necessary-and-sufficient classification of all causal strategies. The exact order characterization below is for a declared family of localized affine-square tests, not for the existence of arbitrary history-dependent balancing programs. These distinctions are part of the mathematical statement, not qualifications added after the proof.
